%%
%% part2-book-components/ch09-math.tex
%%
%% Chapter 9: Mathematics — Equations, matrices, operators,
%% delimiters, and number sets.
%%

\chapter{ریاضیات}
\label{chap:math}

\section{چرا این موضوع مهم است؟}
\label{sec:math-why}

ریاضیات، ستون فقرات کتاب‌های فنی و علمی است.
از معادلات ساده تا ماتریس‌های پیچیده، از
عملگرهای حسابان تا مجموعه‌های عددی، همه در
\lr{MatinBook} پشتیبانی می‌شوند.

\lr{MatinBook} از بسته‌های \lr{amsmath}،
\lr{mathtools} و \lr{unicode-math} استفاده
می‌کند و مجموعه‌ای از عملگرها و جداکننده‌های
سفارشی را نیز اضافه کرده است. در این فصل،
با تمام این قابلیت‌ها آشنا می‌شوید.

\mnote{فونت ریاضی \lr{MatinBook}، \lr{XITS Math} است که از \lr{OpenType} پشتیبانی می‌کند و نمادهای ریاضی را با کیفیت بالا نمایش می‌دهد.}

\section{معادلات درون‌خطی}
\label{sec:math-inline}

معادلات درون‌خطی، برای اشاره به نمادها و
عبارات کوتاه استفاده می‌شوند. در \lr{MatinBook}،
این معادلات باید داخل \cmd{lr} قرار گیرند تا
جهت‌نمایی درست حفظ شود:

\begin{latin}
\begin{minted}{latex}
|\pc{رابطه‌ی معروف اینشتین:}|\lr{$E = mc^{2}$}|\pc{است.}|%
|\pc{قضیه‌ی فیثاغورث:}|\lr{$a^{2} + b^{2} = c^{2}$}|\pc{است.}|%
|\pc{اتحاد اویلر:}|\lr{$e^{i\pi} + 1 = 0$}|\pc{است.}|%
\end{minted}
\end{latin}

خروجی:

رابطه‌ی معروف اینشتین: \lr{$E = mc^{2}$} است.
قضیه‌ی فیثاغورث: \lr{$a^{2} + b^{2} = c^{2}$} است.
اتحاد اویلر: \lr{$e^{i\pi} + 1 = 0$} است.

\mnote{همیشه معادلات درون‌خطی را داخل \lr{\textbackslash lr\{...\}} قرار دهید. این کار از به‌هم‌ریختن جهت متن جلوگیری می‌کند.}

\section{معادلات نمایشی}
\label{sec:math-display}

معادلات نمایشی، در خط جداگانه و وسط‌چین نمایش
داده می‌شوند. \lr{MatinBook} از سه محیط اصلی
پشتیبانی می‌کند:

\subsection{معادله‌ی ساده (\lr{equation})}
\label{sec:math-equation}

برای یک معادله‌ی شماره‌دار، از محیط
\env{equation} استفاده کنید:

\begin{latin}
\begin{minted}{latex}
\begin{equation}
    \int_{0}^{\infty} \frac{x^{3}}{e^{x} - 1} \, dx
    = \frac{\pi^{4}}{15}
    \label{eq:stefan-boltzmann}
\end{equation}
\end{minted}
\end{latin}

خروجی:

\begin{equation}
    \int_{0}^{\infty} \frac{x^{3}}{e^{x} - 1} \, dx
    = \frac{\pi^{4}}{15}
    \label{eq:stefan-boltzmann-ch9}
\end{equation}

\mnote{شماره‌گذاری معادلات در \lr{MatinBook} به‌صورت «شماره-فصل» است: \lr{(۱-۲)} یعنی فرمول اول از فصل دوم.}

\subsection{معادلات هم‌تراز (\lr{align})}
\label{sec:math-align}

برای چند معادله‌ی هم‌تراز، از محیط \env{align}
استفاده کنید:

\begin{latin}
\begin{minted}{latex}
\begin{align}
    (a + b)^{2} &= (a + b)(a + b) \nonumber \\
    &= a^{2} + ab + ba + b^{2} \nonumber \\
    &= a^{2} + 2ab + b^{2} \label{eq:binomial}
\end{align}
\end{minted}
\end{latin}

خروجی:

\begin{align}
    (a + b)^{2} &= (a + b)(a + b) \nonumber \\
    &= a^{2} + ab + ba + b^{2} \nonumber \\
    &= a^{2} + 2ab + b^{2} \label{eq:binomial-ch9}
\end{align}

\mnote{دستور \lr{\textbackslash nonumber} معادله را بدون شماره می‌کند. برای معادلات میانی که به آن‌ها ارجاع نمی‌دهید، از این دستور استفاده کنید.}

\subsection{معادلات غیرهم‌تراز (\lr{gather})}
\label{sec:math-gather}

برای معادلات غیرهم‌تراز، از محیط \env{gather}
استفاده کنید:

\begin{latin}
\begin{minted}{latex}
\begin{gather}
    x = a + b + c + d + e + f \label{eq:long1} \\
    y = \alpha + \beta + \gamma + \delta + \epsilon \label{eq:long2}
\end{gather}
\end{minted}
\end{latin}

خروجی:

\begin{gather}
    x = a + b + c + d + e + f \label{eq:long1-ch9} \\
    y = \alpha + \beta + \gamma + \delta + \epsilon \label{eq:long2-ch9}
\end{gather}

\mnote{تفاوت \lr{align} و \lr{gather}: در \lr{align}، معادلات با \lr{\&} هم‌تراز می‌شوند، ولی در \lr{gather}، هر معادله در خط جداگانه و بدون هم‌ترازی نمایش داده می‌شود.}

\section{ماتریس‌ها}
\label{sec:math-matrices}

\lr{MatinBook} سه دستور برای نمایش ماتریس
ارائه می‌دهد:

\begin{itemize}
    \item \cmd{mat\{...\}}: ماتریس با کروشه.
    \item \cmd{pmat\{...\}}: ماتریس با پرانتز.
    \item \cmd{vmat\{...\}}: ماتریس با خط عمودی
    (دترمینان).
\end{itemize}

\subsection{ماتریس با کروشه (\cmd{mat})}
\label{sec:math-mat-bracket}

\begin{latin}
\begin{minted}{latex}
\[
    A = \mat{1 & 2 & 3 \\ 4 & 5 & 6 \\ 7 & 8 & 9}
\]
\end{minted}
\end{latin}

خروجی:

\[
    A = \mat{1 & 2 & 3 \\ 4 & 5 & 6 \\ 7 & 8 & 9}
\]

\subsection{ماتریس با پرانتز (\cmd{pmat})}
\label{sec:math-mat-paren}

\begin{latin}
\begin{minted}{latex}
\[
    B = \pmat{a_{11} & a_{12} \\ a_{21} & a_{22}}
\]
\end{minted}
\end{latin}

خروجی:

\[
    B = \pmat{a_{11} & a_{12} \\ a_{21} & a_{22}}
\]

\subsection{دترمینان (\cmd{vmat})}
\label{sec:math-mat-vmatrix}

\begin{latin}
\begin{minted}{latex}
\[
    \det(B) = \vmat{a_{11} & a_{12} \\ a_{21} & a_{22}}
    = a_{11}a_{22} - a_{12}a_{21}
\]
\end{minted}
\end{latin}

خروجی:

\[
    \det(B) = \vmat{a_{11} & a_{12} \\ a_{21} & a_{22}}
    = a_{11}a_{22} - a_{12}a_{21}
\]

\mnote{برای ماتریس‌های بزرگ، می‌توانید از \lr{\textbackslash cdots}، \lr{\textbackslash vdots} و \lr{\textbackslash ddots} برای نشان دادن ادامه‌ی الگو استفاده کنید.}

\section{عملگرهای ریاضی}
\label{sec:math-operators}

\lr{MatinBook} مجموعه‌ای از عملگرهای سفارشی
ارائه می‌دهد که در حوزه‌های مختلف ریاضیات
و علوم کامپیوتر کاربرد دارند. جدول
\cref{tab:math-operators} این عملگرها را
خلاصه می‌کند.

\begin{matintable}{عملگرهای ریاضی \lr{MatinBook}}{tab:math-operators}
\begin{tblr}{
    width = \textwidth,
    colspec = {l l X[l]},
    hlines, vlines,
    colsep = 6pt,
    rowsep = 4pt,
    row{1} = {font=\bfseries},
}
دسته & دستور & خروجی \\
حسابان برداری & \lr{\textbackslash grad} & \lr{$\grad f$} \\
حسابان برداری & \lr{\textbackslash curl} & \lr{$\curl \vec{F}$} \\
حسابان برداری & \lr{\textbackslash diver} & \lr{$\diver \vec{F}$} \\
جبر خطی & \lr{\textbackslash rank} & \lr{$\rank(A)$} \\
جبر خطی & \lr{\textbackslash trace} & \lr{$\trace(A)$} \\
جبر خطی & \lr{\textbackslash diag} & \lr{$\diag(a, b, c)$} \\
بهینه‌سازی & \lr{\textbackslash argmax} & \lr{$\argmax_x f(x)$} \\
بهینه‌سازی & \lr{\textbackslash argmin} & \lr{$\argmin_x f(x)$} \\
نظریه اعداد & \lr{\textbackslash gcdop} & \lr{$\gcdop(a, b)$} \\
نظریه اعداد & \lr{\textbackslash lcm} & \lr{$\lcm(a, b)$} \\
توابع خاص & \lr{\textbackslash sinc} & \lr{$\sinc(x)$} \\
توابع خاص & \lr{\textbackslash sgn} & \lr{$\sgn(x)$} \\
توابع خاص & \lr{\textbackslash erf} & \lr{$\erf(x)$} \\
\end{tblr}
\end{matintable}

\mnote{نام \lr{\textbackslash gcdop} از \lr{gcd operator} گرفته شده، چون \lr{\textbackslash gcd} در برخی بسته‌ها تعریف شده است.}

\section{جداکننده‌های هوشمند}
\label{sec:math-delimiters}

\lr{MatinBook} شش جداکننده‌ی هوشمند ارائه می‌دهد
که به‌طور خودکار با محتوای خود مقیاس می‌گیرند.
جدول \cref{tab:math-delimiters} این جداکننده‌ها
را خلاصه می‌کند.

\begin{matintable}{جداکننده‌های هوشمند}{tab:math-delimiters}
\begin{tblr}{
    width = \textwidth,
    colspec = {l l X[l]},
    hlines, vlines,
    colsep = 6pt,
    rowsep = 4pt,
    row{1} = {font=\bfseries},
}
دستور & خروجی & کاربرد \\
\lr{\textbackslash abs\{x\}} & \lr{$\abs{x}$} & قدر مطلق \\
\lr{\textbackslash norm\{x\}} & \lr{$\norm{x}$} & نرم \\
\lr{\textbackslash ceil\{x\}} & \lr{$\ceil{x}$} & سقف \\
\lr{\textbackslash floor\{x\}} & \lr{$\floor{x}$} & کف \\
\lr{\textbackslash inner\{u\}\{v\}} & \lr{$\inner{u}{v}$} & ضرب داخلی \\
\lr{\textbackslash paren\{x\}} & \lr{$\paren{x}$} & پرانتز هوشمند \\
\end{tblr}
\end{matintable}

\subsection{مقیاس خودکار}
\label{sec:math-delimiters-scaling}

این جداکننده‌ها به‌طور خودکار با محتوای خود
مقیاس می‌گیرند:

\begin{latin}
\begin{minted}{latex}
\[
    \abs{\frac{a}{b}}, \quad
    \norm{\sum_{i=1}^{n} x_i}, \quad
    \paren{\frac{a + b}{c + d}}
\]
\end{minted}
\end{latin}

خروجی:

\[
    \abs{\frac{a}{b}}, \quad
    \norm{\sum_{i=1}^{n} x_i}, \quad
    \paren{\frac{a + b}{c + d}}
\]

\mnote{برای مقیاس دستی، می‌توانید از \lr{\textbackslash abs[\textbackslash Big]\{...\}}، \lr{\textbackslash abs[\textbackslash big]\{...\}} و... استفاده کنید.}

\section{مجموعه‌های عددی}
\label{sec:math-numbers}

\lr{MatinBook} پنج مجموعه‌ی عددی استاندارد
ارائه می‌دهد:

\begin{itemize}
    \item \cmd{N}: مجموعه‌ی اعداد طبیعی
    (\lr{$\N$}).
    \item \cmd{Z}: مجموعه‌ی اعداد صحیح
    (\lr{$\Z$}).
    \item \cmd{Q}: مجموعه‌ی اعداد گویا
    (\lr{$\Q$}).
    \item \cmd{R}: مجموعه‌ی اعداد حقیقی
    (\lr{$\R$}).
    \item \cmd{C}: مجموعه‌ی اعداد مختلط
    (\lr{$\C$}).
\end{itemize}

رابطه‌ی شمول میان این مجموعه‌ها:

\begin{latin}
\begin{minted}{latex}
\[
    \N \subset \Z \subset \Q \subset \R \subset \C
\]
\end{minted}
\end{latin}

خروجی:

\[
    \N \subset \Z \subset \Q \subset \R \subset \C
\]

\mnote{این مجموعه‌ها با \lr{\textbackslash providecommand} تعریف شده‌اند، پس می‌توانید آن‌ها را با \lr{\textbackslash renewcommand} بازنویسی کنید.}

\section{نمادگذاری مجموعه}
\label{sec:math-sets}

\lr{MatinBook} سه دستور برای نمادگذاری مجموعه
ارائه می‌دهد:

\begin{itemize}
    \item \cmd{set\{...\}}: مجموعه‌ی ساده.
    \item \cmd{setbuilder\{...\}\{...\}}: مجموعه‌ی
    سازنده.
    \item \cmd{given}: جداکننده‌ی شرط.
\end{itemize}

\begin{latin}
\begin{minted}{latex}
\[
    A = \set{1, 2, 3, 4, 5}, \quad
    B = \setbuilder{x \in \R}{x > 0}
\]
\end{minted}
\end{latin}

خروجی:

\[
    A = \set{1, 2, 3, 4, 5}, \quad
    B = \setbuilder{x \in \R}{x > 0}
\]

\mnote{دستور \lr{\textbackslash setbuilder} دو آرگومان می‌گیرد: شرط و محدوده.}

\section{مشتق و انتگرال}
\label{sec:math-derivatives}

\lr{MatinBook} چهار دستور برای مشتق و انتگرال
ارائه می‌دهد:

\begin{itemize}
    \item \cmd{dd\{x\}}: دیفرانسیل (\lr{$dx$}).
    \item \cmd{D}: دیفرانسیل کوتاه (\lr{$\D$}).
    \item \cmd{deriv\{f\}\{x\}}: مشتق معمولی
    (\lr{$\deriv{f}{x}$}).
    \item \cmd{pderiv\{f\}\{x\}}: مشتق جزئی
    (\lr{$\pderiv{f}{x}$}).
\end{itemize}

\begin{latin}
\begin{minted}{latex}
\[
    \deriv{f}{x} = \lim_{h \to 0} \frac{f(x + h) - f(x)}{h}, \quad
    \pderiv{f}{x} = \frac{\partial f}{\partial x}
\]
\end{minted}
\end{latin}

خروجی:

\[
    \deriv{f}{x} = \lim_{h \to 0} \frac{f(x + h) - f(x)}{h}, \quad
    \pderiv{f}{x} = \frac{\partial f}{\partial x}
\]

\mnote{استفاده از \lr{\textbackslash dd\{x\}} به‌جای \lr{dx} در انتگرال‌ها، فاصله‌ی بهتری ایجاد می‌کند.}

\section{شماره‌گذاری معادلات}
\label{sec:math-equation-numbering}

شماره‌گذاری معادلات در \lr{MatinBook} به‌صورت
\textbf{«شماره-فصل»} است. یعنی معادله‌ی شماره‌ی
\lr{(۵-۱)} به معنای فرمول پنجم از فصل اول است.

\begin{matintable}{نمونه‌ی شماره‌گذاری معادلات}{tab:equation-numbering}
\begin{tblr}{
    width = 0.7\textwidth,
    colspec = {c c c},
    hlines, vlines,
    colsep = 8pt,
    rowsep = 4pt,
    row{1} = {font=\bfseries},
}
فصل & معادله & نمایش \\
۱ & اول & \lr{(۱-۱)} \\
۱ & دوم & \lr{(۲-۱)} \\
۲ & اول & \lr{(۱-۲)} \\
۲ & دوم & \lr{(۲-۲)} \\
\end{tblr}
\end{matintable}

\mnote{این شماره‌گذاری، مطابق استاندارد نشر دانشگاهی ایران (دانشگاه تهران و دانشگاه علوم پزشکی) است.}

\section{شکستن معادلات طولانی}
\label{sec:math-breaking}

\lr{MatinBook} امکان شکستن معادلات طولانی بین
صفحات را فراهم می‌کند:

\begin{latin}
\begin{minted}{latex}
\begin{align}
    f(x) &= a_{0} + a_{1}x + a_{2}x^{2} + \cdots \nonumber \\
         &\quad + a_{n}x^{n} \\
    &= \sum_{i=0}^{n} a_{i}x^{i}
\end{align}
\end{minted}
\end{latin}

خروجی:

\begin{align}
    f(x) &= a_{0} + a_{1}x + a_{2}x^{2} + \cdots \nonumber \\
         &\quad + a_{n}x^{n} \\
    &= \sum_{i=0}^{n} a_{i}x^{i}
\end{align}

\mnote{دستور \lr{\textbackslash allowdisplaybreaks} در \lr{MatinBook} فعال است، پس معادلات طولانی به‌طور خودکار بین صفحات می‌شکنند.}

\section{خطاهای رایج}
\label{sec:math-mistakes}

\subsection{فراموش کردن \cmd{lr} برای معادلات درون‌خطی}
\label{sec:math-mistake-lr}

اگر معادله‌ی درون‌خطی را داخل \cmd{lr} قرار
ندهید، ممکن است جهت آن به‌هم بریزد:

\begin{latin}
\begin{minted}{latex}
% |\pc{اشتباه}|:
|\pc{این معادله}|\lr{$E = mc^{2}$}|\pc{است.}|%

% |\pc{درست}|:
|\pc{این معادله}|\lr{$E = mc^{2}$}|\pc{است.}|%
\end{minted}
\end{latin}

\subsection{استفاده از \cmd{\$...\$} در عنوان بخش}
\label{sec:math-mistake-title}

هرگز از حالت ریاضی در عنوان بخش استفاده نکنید:

\begin{latin}
\begin{minted}{latex}
% |\pc{اشتباه}|:
\section{|\pc{تحلیل}|\lr{$O(n)$}}

% |\pc{درست}|:
\section{|\pc{تحلیل}|\lr{O(n)}}
\end{minted}
\end{latin}

\subsection{فراموش کردن شماره‌گذاری معادلات مهم}
\label{sec:math-mistake-label}

معادلاتی که به آن‌ها ارجاع می‌دهید، باید
شماره‌دار باشند:

\begin{latin}
\begin{minted}{latex}
% |\pc{اشتباه}|:
\[
    E = mc^{2}
\]
|\pc{طبق معادله‌ی بالا...}|%

% |\pc{درست}|:
\begin{equation}
    E = mc^{2}
    \label{eq:einstein}
\end{equation}
|\pc{طبق}|\cref{eq:einstein}|\pc{...}|%
\end{minted}
\end{latin}

\section{تمرین‌ها}
\label{sec:math-exercises}

\begin{exercise}
    یک معادله‌ی درون‌خطی و یک معادله‌ی نمایشی
    بنویسید. به معادله‌ی نمایشی ارجاع دهید.
    \label{exc:math-equation}
\end{exercise}

\begin{exercise}
    یک ماتریس ۳×۳ با کروشه و یک ماتریس ۲×۲
    با پرانتز بنویسید.
    \label{exc:math-matrix}
\end{exercise}

\begin{exercise}
    از سه عملگر سفارشی \lr{MatinBook} در یک
    معادله استفاده کنید.
    \label{exc:math-operators}
\end{exercise}

\begin{exercise}
    یک مجموعه‌ی سازنده بنویسید که اعداد
    حقیقی مثبت را توصیف کند.
    \label{exc:math-set}
\end{exercise}

\section{خلاصه}
\label{sec:math-summary}

در این فصل، با ریاضیات در \lr{MatinBook}
آشنا شدیم:

\begin{itemize}
    \item معادلات درون‌خطی باید داخل \cmd{lr}
    قرار گیرند.

    \item معادلات نمایشی با سه محیط
    \env{equation}، \env{align} و \env{gather}
    نمایش داده می‌شوند.

    \item سه دستور برای ماتریس:
    \cmd{mat}، \cmd{pmat}، \cmd{vmat}.

    \item ۱۳ عملگر سفارشی برای حسابان برداری،
    جبر خطی، بهینه‌سازی، نظریه اعداد و توابع
    خاص.

    \item شش جداکننده‌ی هوشمند که با محتوا
    مقیاس می‌گیرند.

    \item پنج مجموعه‌ی عددی (\lr{$\N$}،
    \lr{$\Z$}، \lr{$\Q$}، \lr{$\R$}، \lr{$\C$}).

    \item شماره‌گذاری معادلات به‌صورت
    «شماره-فصل» است.

    \item خطاهای رایج شامل فراموش کردن
    \cmd{lr}، استفاده از حالت ریاضی در
    عنوان بخش، و فراموش کردن شماره‌گذاری
    معادلات مهم است.
\end{itemize}

در فصل بعدی (\cref{chap:code})، با کد و
الگوریتم در \lr{MatinBook} آشنا می‌شوید و
یاد می‌گیرید که چگونه کدهای برنامه‌نویسی
و شبه‌کدها را نمایش دهید.